% =============================================================================
% Chapter 8 -- Experimental Setup
% =============================================================================

\chapter{Experimental Setup}
\label{ch:experimental_setup}

This chapter details the experimental methodology used to evaluate the proposed circular oversampling algorithms.
We describe the benchmark datasets, classifiers, evaluation metrics, validation protocol, comparison methods, hyperparameter settings, and statistical testing framework.


% ============================================================================
\section{Benchmark Datasets}
\label{sec:setup:datasets}

We evaluate the proposed methods on 14 binary classification datasets from the KEEL~\cite{alcala2011keel} and UCI~\cite{uci2017} repositories.
These datasets are widely used in the imbalanced learning literature and span a range of imbalance ratios, dimensionalities, and sample sizes.

\begin{table}[H]
\centering
\caption{Summary of the 14 benchmark datasets. IR = Imbalance Ratio ($n_{\text{maj}}/n_{\text{min}}$).}
\label{tab:setup:datasets}
\begin{tabular}{llccccl}
\toprule
\textbf{No.} & \textbf{Dataset} & \textbf{Source} & \textbf{Samples} & \textbf{Features} & \textbf{IR} & \textbf{Category} \\
\midrule
1  & heart         & UCI  & 303  & 13 & 1.25  & Low \\
2  & ionosphere    & UCI  & 351  & 34 & 1.79  & Low \\
3  & glass1        & KEEL & 214  & 9  & 1.82  & Low \\
4  & wisconsin     & UCI  & 699  & 9  & 1.86  & Low \\
5  & pima          & UCI  & 768  & 8  & 1.87  & Low \\
6  & yeast1        & KEEL & 1484 & 8  & 2.46  & Moderate \\
7  & haberman      & UCI  & 306  & 3  & 2.78  & Moderate \\
8  & vehicle0      & KEEL & 846  & 18 & 3.25  & Moderate \\
9  & ecoli1        & KEEL & 336  & 7  & 3.36  & Moderate \\
10 & new-thyroid1  & KEEL & 215  & 5  & 5.14  & High \\
11 & ecoli2        & KEEL & 336  & 7  & 5.46  & High \\
12 & yeast3        & KEEL & 1484 & 8  & 8.10  & High \\
13 & ecoli3        & KEEL & 336  & 7  & 8.60  & High \\
14 & glass4        & KEEL & 214  & 9  & 15.47 & Extreme \\
\bottomrule
\end{tabular}
\end{table}

The datasets are grouped into four IR categories:
\begin{itemize}[leftmargin=*]
    \item \textbf{Low IR} (1.0--2.0): heart, ionosphere, glass1, wisconsin, pima.
    \item \textbf{Moderate IR} (2.0--5.0): yeast1, haberman, vehicle0, ecoli1.
    \item \textbf{High IR} (5.0--10.0): new-thyroid1, ecoli2, yeast3, ecoli3.
    \item \textbf{Extreme IR} ($>$10.0): glass4.
\end{itemize}


\subsection{PCA Explained Variance}
\label{sec:setup:pca_variance}

Since all three proposed algorithms project minority instances to a 2-D PCA subspace before generating synthetic samples (Section~\ref{sec:methods:framework}), the fraction of variance captured by the two principal components determines how faithfully the circular geometry in 2-D approximates the actual minority-class structure in the original feature space.

\begin{table}[H]
\centering
\caption{Explained variance ratio of the first two principal components of the minority class per dataset.
Low values indicate that the 2-D approximation captures a smaller fraction of the original geometry.}
\label{tab:setup:pca_variance}
\begin{tabular}{lcccc}
\toprule
\textbf{Dataset} & \textbf{Features} & \textbf{PC1} & \textbf{PC2} & \textbf{PC1+PC2} \\
\midrule
ecoli1       & 7  & 0.82 & 0.12 & 0.94 \\
ecoli2       & 7  & 0.79 & 0.13 & 0.92 \\
ecoli3       & 7  & 0.77 & 0.14 & 0.91 \\
glass1       & 9  & 0.71 & 0.15 & 0.86 \\
glass4       & 9  & 0.68 & 0.17 & 0.85 \\
yeast1       & 8  & 0.76 & 0.14 & 0.90 \\
yeast3       & 8  & 0.73 & 0.16 & 0.89 \\
new-thyroid1 & 5  & 0.91 & 0.06 & 0.97 \\
haberman     & 3  & 0.94 & 0.04 & 0.98 \\
vehicle0     & 18 & 0.59 & 0.14 & 0.73 \\
pima         & 8  & 0.78 & 0.13 & 0.91 \\
wisconsin    & 9  & 0.84 & 0.09 & 0.93 \\
heart        & 13 & 0.62 & 0.12 & 0.74 \\
ionosphere   & 34 & 0.48 & 0.17 & 0.65 \\
\midrule
\textbf{Average} & & 0.74 & 0.13 & 0.87 \\
\bottomrule
\end{tabular}
\end{table}

The 2-D PCA projection retains on average 87\% of the minority-class variance.
For low-dimensional datasets (haberman 3 features, new-thyroid1 5 features), the projection captures 97--98\% of variance and the circular approximation is nearly exact.
For high-dimensional datasets (ionosphere 34 features: 65\%, vehicle0 18 features: 73\%, heart 13 features: 74\%), the circular geometry is a rougher approximation; the inverse PCA transform produces hyper-ellipsoidal generation regions in the original space (Section~\ref{sec:methods:framework}).

\textbf{Implication for native-dimension mode.}
For ionosphere (35\% variance discarded by PCA) and vehicle0 (27\% discarded), the \texttt{use\_pca=False} native-dimension mode is particularly relevant: it performs all geometric operations in the full feature space without discarding minority-class variance.
The ablation in Chapter~\ref{ch:ablation} suggests that native-dimension mode improves performance across the 14 evaluated datasets, with larger gains on higher-dimensional datasets.


% ============================================================================
\section{Classifiers}
\label{sec:setup:classifiers}

Seven classifiers are evaluated: \textbf{KNN} ($k=5$, Euclidean), \textbf{DT} (Gini, unbounded depth), \textbf{SVM-RBF} ($C=1.0$, $\gamma=\text{scale}$), \textbf{NB} (Gaussian), \textbf{MLP} (100 hidden units, ReLU, Adam optimizer~\cite{buda2018systematic}), \textbf{LR} ($L_2$, $C=1.0$), and \textbf{RF} (100 trees, Gini)~\cite{breiman2001random}.

All classifiers are implemented using scikit-learn~\cite{pedregosa2011scikit} with default parameters unless otherwise stated.


% ============================================================================
\section{Evaluation Metrics}
\label{sec:setup:metrics}

Following established practice in the imbalanced learning literature~\cite{he2009learning, branco2016survey, fernandez2018learning}, we report the six metrics defined in Table~\ref{tab:bg:metrics}: F1-score, G-mean, AUC-ROC, Balanced Accuracy, Precision, and Sensitivity.

F1-score, G-mean, and AUC are the primary metrics for method comparison, while balanced accuracy, precision, and sensitivity provide complementary perspectives.


% ============================================================================
\section{Validation Protocol}
\label{sec:setup:validation}

We use 5-fold stratified cross-validation following the oversampling-inside-folds protocol described in Section~\ref{sec:bg:clustering}:

\begin{enumerate}[leftmargin=*]
    \item The dataset is partitioned into 5 stratified folds.
    \item For each fold:
    \begin{enumerate}
        \item Features are standardized (zero mean, unit variance) using training-fold statistics only.
        \item Oversampling is applied to the training fold only.
        \item The classifier is trained on the oversampled training data.
        \item Performance is evaluated on the untouched test fold.
    \end{enumerate}
    \item Results are averaged across all 5 folds.
\end{enumerate}

To reduce variance in the estimates, the entire 5-fold procedure is repeated 3 times with different random seeds, for a total of 15 evaluations per dataset-method-classifier combination.
The reported results are the mean across all 15 evaluations.


% ============================================================================
\section{Comparison Methods}
\label{sec:setup:comparison}

The proposed methods are compared against the following baselines:

\begin{enumerate}[leftmargin=*]
    \item \textbf{None}: No oversampling (original imbalanced data).
    \item \textbf{Random Oversampling (ROS)}~\cite{batista2004study}: Duplicates random minority instances.
    \item \textbf{SMOTE}~\cite{chawla2002smote}: Standard linear interpolation ($k=5$).
    \item \textbf{Borderline-SMOTE}~\cite{han2005borderline}: Borderline-focused SMOTE (variant 1, $k=5$).
    \item \textbf{ADASYN}~\cite{he2008adasyn}: Adaptive synthetic oversampling ($k=5$).
    \item \textbf{SVM-SMOTE}~\cite{nguyen2011borderline}: SVM-based borderline oversampling ($k=5$).
    \item \textbf{K-Means SMOTE}~\cite{douzas2018improving}: Clustering-based SMOTE (K determined by silhouette score).
    \item \textbf{Circular-SMOTE}: Our baseline circular oversampler with uniform disk sampling (Section~\ref{sec:methods:framework}).
\end{enumerate}

Safe-Level SMOTE~\cite{bunkhumpornpat2009safe} was reviewed in Chapter~\ref{ch:literature} but excluded from the main comparison: the \texttt{imbalanced-learn} implementation raises an exception when the safe level equals zero for all minority instances in a fold, which occurs on 3 of the 14 benchmark datasets (glass4, new-thyroid1, ecoli3) due to extreme imbalance.

Baselines 2--7 are implemented using the imbalanced-learn library~\cite{lemaitre2017imbalanced}.
Circular-SMOTE and the proposed methods are implemented in a custom framework.


% ============================================================================
\section{Hyperparameter Settings}
\label{sec:setup:hyperparams}

Table~\ref{tab:setup:hyperparams} lists the default hyperparameter settings for the proposed methods.

\begin{table}[H]
\centering
\caption{Default hyperparameters for the proposed methods.}
\label{tab:setup:hyperparams}
\begin{tabular}{llll}
\toprule
\textbf{Parameter} & \textbf{Symbol} & \textbf{Default} & \textbf{Used in} \\
\midrule
Number of clusters        & $K$               & 3    & All \\
Seed neighbors           & $k_{\text{seed}}$  & 5    & All \\
Local density neighbors  & $k_{\text{nn}}$   & 5    & GVM-CO, LS-CO \\
Max concentration         & $\kappa_{\max}$   & 20   & GVM-CO, LS-CO \\
Concentration exponent    & $\gamma$          & 2.0  & GVM-CO, LRE-CO \\
Score blending            & $\alpha$          & 0.6  & GVM-CO, LRE-CO \\
Certainty threshold       & $\tau$            & 0.80 & LRE-CO \\
Clustering method         & --                & K-Means & All \\
Denoise method            & --                & None & All \\
\midrule
\multicolumn{4}{l}{\textit{LS-CO Specific}} \\
\midrule
Number of layers          & $L$               & 10$^*$ & LS-CO \\
Number of segments        & $M$               & 8    & LS-CO \\
Layer decay               & $\beta$           & 1.0  & LS-CO \\
Angular decay             & $\sigma_\theta$   & $\pi/4$ & LS-CO \\
\midrule
\multicolumn{4}{l}{\textit{Seed Selection}} \\
\midrule
Candidates                & $N_{\text{cand}}$  & 100  & -- \\
PCA components            & $d_{\text{PCA}}$  & 2    & -- \\
Clusters for seeds        & $K_s$             & 5    & -- \\
Topological $k$           & $k_{\text{topo}}$ & 5    & -- \\
Smoothness penalty        & $w_z$             & 0.5  & -- \\
\bottomrule
\end{tabular}
\end{table}

For all methods, the sampling strategy achieves a 1:1 minority-to-majority ratio after oversampling.

\noindent$^*$\textit{The main comparison uses $L=10$ as a balanced default.
The ablation in Chapter~\ref{ch:ablation} shows that $L \geq 60$ achieves higher F1 and G-Mean while remaining under 1~s/fold; practitioners without strict runtime constraints should use $L=60$.}

\medskip
\noindent\textbf{Hyperparameter selection and optimistic bias.}
The default hyperparameter values listed above (including $\tau=0.80$, $\kappa_{\max}=20$, $L=10$) were selected based on ablation experiments conducted on the same 14 benchmark datasets used for the main comparison.
This constitutes a form of indirect leakage: the defaults are empirically best-performing values on the final evaluation benchmark.
The expected optimistic bias has not been formally quantified via nested cross-validation; practitioners applying these defaults to new datasets should be aware that performance may differ.


% ============================================================================
\section{Statistical Testing Framework}
\label{sec:setup:stats}

To assess the statistical significance of performance differences, we follow the recommendations of Dem\v{s}ar~\cite{demsar2006statistical} and Garc\'{i}a and Herrera~\cite{garcia2008extension}:

\begin{enumerate}[leftmargin=*]
    \item \textbf{Friedman test}: Tests the null hypothesis that all methods have equal average ranks across datasets ($\alpha = 0.05$).
    \item \textbf{Holm's step-down correction}: Applied to pairwise comparisons when the Friedman test rejects the null.
    \item \textbf{Nemenyi test}: Computes the critical difference for a CD diagram ($\alpha = 0.05$).
\end{enumerate}

All statistical tests are implemented using SciPy~\cite{pedregosa2011scikit} and scikit-posthocs.


% ============================================================================
\section{Implementation Details}
\label{sec:setup:implementation}

The entire experimental framework is implemented in Python 3.10.
The core oversampling algorithms use NumPy for vectorised computation, scikit-learn for clustering, KNN, PCA, and classifiers, and imbalanced-learn for baseline methods.

All experiments are run on a machine with an Intel i7-12700H processor and 32 GB RAM.
No GPU acceleration is used.
Random seeds are fixed for reproducibility: the global seed is set to 42, and cross-validation seeds are $\{42, 123, 456\}$ for the three repetitions.


